\section{Conclusion and future work}
\label{sec:conclusion}

\subsection{What was set out to do, and what was delivered}

This thesis set out to apply reinforcement learning to a custom high-speed biped. The platform was
inherited as a partially finished mechanical design. Three things were delivered.

A \textbf{working method}. It comprises the modelling, training and curriculum stack that produces
a 2.6\,m/s never-falling runner on the free simulated plant. It also comprises the CAD-to-MJCF
pipeline that closes a kinematic loop URDF cannot express. And it comprises the deployment stack
that closes a 200\,Hz loop on a non-realtime single-board computer.

A \textbf{measured robot}. Every load-bearing parameter of DASH-01 is replaced by a measurement,
with stated uncertainty and stated un-resolvables. The identification reports a repeatable
per-joint gravity gain. It also reports explicitly that the mass distribution is not the
explanation for the residual, so the next person does not fit the same null space again.

An \textbf{account of the distance between the two}. The measured robot is 18\,\% heavier and
15\,\% weaker than the model it trained against, and an order of magnitude slower in actuation. Its
corrected ceiling is about 2.7\,m/s burst and 1.3\,m/s sustained. That reclassifies the platform
from record contender to demonstrator, and \cref{sec:future} ranks the changes that would move it.

\subsection{The answer to the research question}

\Cref{sec:bg-gap} asked at what level learning should sit, given this robot's actuator dynamics and
available control bandwidth. The answer this work supports is narrow. \textbf{On a plant whose
actuation bandwidth is an order of magnitude below its own divergence timescale, the level at which
learning sits matters less than the level at which the actuator is commanded.}

The evidence is that a learned controller and a 21-gain scripted controller hit the identical pitch
wall at the identical time. The wall is predicted by a formula containing no controller term. The
learning method held wherever it could be tested, and classical control matched it at the plant's
own limits. That agreement is a consistency check on the stack. It is not evidence that either
method won, because the plant bound both.

\subsection{Limitations}

Recorded throughout, and gathered once here. No policy has run on the hardware, and export plus
inference on the Raspberry Pi remain the critical path. The sprint results rest on privileged
observations, namely base linear velocity and distance to the line. Identification left the link
inertia tensors and the 0.9\,N$\cdot$m gravity residual unresolved, with the kinematic model as the
remaining suspect. The actuator model lacks a coupled torque--speed envelope, which makes the
simulation optimistic exactly in the high-speed regime the project cares about. The corrected
capability figures propagate sensitivity coefficients rather than re-running the feasibility sweep.
The impedance action space was designed and never trained. And the physical robot currently stands
with one functional leg.

\subsection{Future work}

For \emph{this} robot the record orders five steps, and \cref{sec:future} gives them with their
dependencies. First, ONNX export and a Pi inference loop. Second, current-mode control with the
position PD closed on the Pi. Third, the carbon ankle retrain. Fourth, a from-scratch retrain with
cadence pressure from step zero. Fifth, the phase-dependent impedance experiment that the first two
make possible.

For a \emph{successor} robot the ranking is bus voltage, then distal mass, then the 4-bar
proportions, then ankle sizing. Motor cooling is a lever with a named test, and torque-mode control
is a design opinion. Underneath all of them sits one method-level recommendation. A validated force
map and a validated feasibility sweep now exist. The natural next use of the training cluster is
therefore a design search over robots, rather than a further search over controllers for a robot
already known to be the constraint.
